\clearpage
\section*{Sources and further reading}
\addcontentsline{toc}{section}{Sources and further reading}
\subsection*{Project sources}
\begin{enumerate}[label={[P\arabic*]},leftmargin=2.5em]
\item Komorebi P. \emph{Projector Refinement Physics, Volume I}. The foundational source studies cover transfer calculus, mixed and open dynamics, countable resolutions, unbounded domains, positive effects, and matrix-corner derivations. The technical source concordance identifies the relevant statements.
\item Komorebi P. \emph{Projector Refinement Physics, Volume I: Technical Supplement}. Claim statements, formula registries, computational checks, and the hypotheses delimiting their conclusions.
\item Komorebi P. \emph{Komorebi Physics formal source}. The foundations verification appendix gives declaration identities, retained compiler evidence, and the limits of its coverage.
\item Komorebi P. \emph{Fourier Packets and the Continuum Dirac Operator}. The free-field construction in this book.
\end{enumerate}

\subsection*{Classical context}
\begin{enumerate}[label={[C\arabic*]},leftmargin=2.5em]
\item G. Lindblad. ``On the generators of quantum dynamical semigroups.'' \emph{Communications in Mathematical Physics} \textbf{48} (1976), 119--130. \url{https://doi.org/10.1007/BF01608499}.
\item V. Gorini, A. Kossakowski, and E. C. G. Sudarshan. ``Completely positive dynamical semigroups of N-level systems.'' \emph{Journal of Mathematical Physics} \textbf{17}, no.5 (1976), 821--825. \url{https://doi.org/10.1063/1.522979}.
\item John Watrous. \emph{Theory of Quantum Information}, CS766/QIC820, Fall2011, Lecture5, Section5.1, Theorem5.1. The finite-dimensional Naimark construction is given by the square-root isometry and coordinate projectors. \url{https://cs.uwaterloo.ca/~watrous/TQI-notes/TQI-notes.05.pdf}.
\end{enumerate}

The first two classical references identify the standard setting for the finite Lindblad expression used in this volume. Their general classification results are not premises of the elementary packet calculations proved here. The Watrous notes give a direct comparison for the finite dilation construction. The source-check record distinguishes verified bibliographic metadata, the inspected dilation proof, and the inaccessible full AIP article.

This bibliography provides context and traceability. It is not an exhaustive literature review or a claim of historical priority. The mathematical review concerns the displayed arguments and the explicitly recorded source correspondence.
